\section{Repair Tests and Training Details}
\label{app:protocol}
This appendix specifies the state restoration, evidence handling, and learning
procedure used in Section~\ref{sec:method}.

\subsection{Notation and Information Available to Each Role}
\label{app:views}
\begin{table}[htbp]
\caption{Notation for a repair, an improvement round, and an evaluation task.}
\centering
\begin{tabularx}{\linewidth}{lX}
\toprule
Symbol & Meaning\\\midrule
$\pi_k$ & Retained task policy entering round $k$.\\
$\widetilde\pi_{k+1}$ & Trained candidate awaiting the selection test.\\
$x,\ c(x)$ & Recoverable environment state and policy-visible context.\\
$r,\ a_0$ & Registered repair and current-policy baseline action.\\
$E_k,\ S_k$ & Current trajectory evidence and round-start experience snapshot.\\
$\mathcal U$ & Common predeclared set of failure units.\\
$\mathcal I_k$ & Completed interventions whose source and training records match the executed test.\\
$J_{\mathrm{eval}}$ & Repetitions per evaluation task; one in the primary benchmark.\\
$\mathcal D_k^+$ & Benefit subset used for positive supervision.\\
$A_k$ & Scaffold-free acceptance result for candidate policy $k+1$.\\
\bottomrule
\end{tabularx}
\end{table}

\begin{table}[htbp]
\caption{Research inputs and responsibilities. Current experiment outcomes
enter POST directly; historical retrieval reads the round-start snapshot.}
\centering
\begin{tabularx}{\linewidth}{L{0.82in}>{\raggedright\arraybackslash}X>{\raggedright\arraybackslash}X}
\toprule
Role & Available evidence & Responsibility\\\midrule
Task policy & Goal, visible prefix, menu; permitted procedural cues when memory is enabled & Select the next action. OFF/OFF evaluation removes the extra cues.\\
Analyzer & Trajectories, mechanical facts, failure analogies, counterexamples & Propose source-grounded failure hypotheses and repairs.\\
Planner PRE & Failure profiles, candidates, resources, earlier experiments & Select and register the intervention portfolio before current verdicts.\\
Planner POST & Registered PRE, fixed current verdicts and costs, earlier lessons & Interpret experiments and recommend TRAIN or NO-TRAIN.\\
Verifier & Executed paired traces and the fixed effect rule & Assign Benefit, Harm, neutral, or uncertain labels.\\
Acceptance rule & Candidate and parent results on training-side selection tasks & Apply the predeclared promotion and regression criteria.\\
\bottomrule
\end{tabularx}
\end{table}

A historical record retains starting conditions, the action/feedback sequence,
observed state changes, a wrong branch and its consequence, the repair,
applicability conditions, counterexamples, and effect provenance. Retrieval
identifies whether a record is an exact past observation or an analogy. Its
usefulness for the current policy is determined by a new experiment.
Within round $k$, historical reads use the immutable snapshot $S_k$.
New records accumulate in a shadow writeback store and become available as
history through $S_{k+1}$ after round closure. POST receives current verdicts
through its explicit experiment-result input. Hidden restoration state and
future outcomes are excluded from task-policy inputs.

\subsection{Restoring the Source and Executing the Repair}
\label{app:restoration}
The source record identifies the task/gamefile, source prefix and state,
policy checkpoint, prompt, initial menu, memory snapshot, candidate payload,
and remaining budget. Restoration uses an environment-supported snapshot
or a checked replay of the recorded prefix. The experiment records the
backend, replay checks, paired seed protocol, and restored-context equality;
the current validation uses checked public-prefix replay. Its five
continuation seed slots are 17, 31, 47, 73, and 101. All 50 planned
branches of the executable-strategy validation are scientifically complete. The policy receives only
$c(x)$, even when restoration requires additional simulator state.

A one-action repair replaces the registered baseline decision once. Earlier options execute one to four registered actions. In the later
validation, ordered guarded rules select admissible actions using the
current public observation, carried-object evidence, and menu. Each rule
has a finite use count and the program obeys the remaining episode budget.
The executor checks each action against the newly observed menu, debits
actual calls and steps, and records completion or invalidity. If the task is unfinished after the program, $\pi_k$ continues from the
resulting state and remaining budget. The later validation appends the
registered strategy to F1 continuation prompts. F0 retains the original
continuation input. This cue is part of the treatment; it is not added to
the scaffold-free selection prompts.
The intervention remains fixed throughout the paired repetitions.
This tests the whole option, including its extra actions and observations.
The completed training run uses one initial-action example per verified
source, not an expansion of every executed program step. A future
action-only comparison must match this action coverage explicitly.

Missing restoration, unsupported options, parser violations, and incomplete
branches have explicit execution outcomes. The discovery export retains
attempted, supported, completed, and labeled counts over the registered
universe. Only complete protocol-valid tests enter the effect-label rule.
This preserves the distinction between an experiment that found no benefit
and one that did not produce an interpretable comparison.

\subsection{Paired Outcomes and the Stability Rule}
\label{app:labels}
For five valid paired repetitions, define
\begin{equation}
n_{ab}=\sum_{j=1}^{5}\ind\{Y_j^{(0)}=a,Y_j^{(1)}=b\},
\qquad a,b\in\{0,1\}.
\label{eq:counts}
\end{equation}
The effect label is
\begin{equation}
L(x,r)=\begin{cases}
\Benefit\ \text{(Benefit)},&n_{01}\ge4,\ n_{10}=0,\\
\Harm\ \text{(Harm)},&n_{10}\ge4,\ n_{01}=0,\\
\NeutralS\ \text{(neutral-success)},&n_{11}\ge4,\\
\NeutralF\ \text{(neutral-failure)},&n_{00}\ge4,\\
\Uncertain\ \text{(Uncertain)},&\text{otherwise}.
\end{cases}
\label{eq:labels}
\end{equation}
Insufficient valid pairs retain their incomplete-test status. The endpoint
is terminal task success, fixed when PRE registers the experiment.
Intermediate progress, loop reduction, and environment steps support the
behavioral analysis without replacing that endpoint.

The four-of-five rule describes finite-repeat stability, not a confidence
level. Its effect is conditional on the source state and continuation
policy. Because the source episodes were failures, the sampled states may
provide few successful baseline continuations that a repair could damage.
Harm counts are consequently reported for this selected population.
Verification establishes the utility of the executed intervention; claims
about a particular explanatory field require their own interventions.

\subsection{Training Views and Native Supervision}
\label{app:training-details}
The execution view retains the exact I1 policy prompt and one string-valued
\texttt{action} field. The auxiliary view uses the same policy-visible
context followed by a fixed training-only instruction requesting a compact
strategy response. The target describes the principal bottleneck, current
subgoal, expected next event and state change, progress criterion, recovery
trigger, fallback condition, and action. A strategy identifier and evidence
references link this response to the tested repair. The strategy identifier
is computed from the pre-verification payload.
Verification provenance links that same payload to the effect label, which
controls training eligibility. Expected progress is prospective information;
the target excludes realized future outcomes, rewards, Benefit tags,
promotion decisions, and sealed benchmark feedback. The native census contains four execution views and four auxiliary views:
51 action and 2,546 strategy loss-bearing tokens per pass, with 4,096
prompt tokens masked across the eight examples. No prompt truncation,
retokenization, or repeated chat templating is applied.

Both views share one real source state. A common source-state identifier
links them to that state, while a distinct example identifier distinguishes
each view. Native labels mask prompt and padding tokens with $-100$ and
supervise the assistant suffix. The label
census checks each verified pair for nonzero action supervision and nonzero
substantive strategy supervision. It reports action, strategy, and metadata
tokens separately, including the renderer's treatment of strategy identifiers and
evidence references. Every admitted pair is Benefit-verified; the census
rejects a purported strategy view that contains only an action target.
These checks concern explicit supervision. Prompt content can still affect
the conditional distribution learned from assistant targets.

For view $v\in\{e,s\}$, assistant-token mask $m_{ivt}$, and minibatch
$\mathcal B$, the common objective is
\begin{equation}
\mathcal L_{\mathrm{policy}}(\theta;\mathcal B)
=-\frac{1}{Z_{\mathcal B}}
\sum_{(i,v)\in\mathcal B}\sum_t
w_{ivt}m_{ivt}\log\pi_\theta(y_{ivt}\mid p_i^v,y_{iv,<t}).
\label{eq:loss}
\end{equation}
Action-only training samples execution views; dual-view training samples
both views. In the executed run, $w_{ivt}=1$ on supervised assistant tokens and
$Z_{\mathcal B}$ is the number of shifted, unmasked targets across the
gradient-accumulation group. Each microbatch has one example; gradients
from four microbatches accumulate using that common token denominator.
There is no separate equal weighting of the two views. Every one of the
eight rows appears once per pass in the registered deterministic order;
four passes produce eight optimizer updates on one GPU. Thus the longer
strategy responses contribute many more supervised token terms, although
token counts alone do not determine their gradient magnitudes.

The executed update continues the retained parent's trainable LoRA
adapter with a fresh optimizer and scheduler. Its four Benefit sources
are drawn from R5 evidence; the candidate is not initialized from a new
untrained base model. An action-only or equal-compute descendant was not
trained in this validation. Consequently, the result does not estimate
the causal contribution of strategy supervision.

\subsection{Additional Training Conditions}
\label{app:optional}
The historical design also identifies T5 for Benefit--Harm preference
learning and T6 for audited RL. They are optional extensions beyond the
primary supervised comparison. A Benefit pair prefers the repair to its
baseline; a Harm pair reverses that preference. Neutral outcomes alone
supply no such ordering. An executed T5 condition would specify its reference
policy, preference serialization, and objective parameters, including
$\beta$ for DPO \citep{rafailov2023dpo}. T6 requires its own executed loss
and rollout protocol. Neither extension contributes a result to this study.

\subsection{Recording Autonomy and Recovery}
Routine scientific decisions include choosing a research question and test
portfolio, recommending training, and accepting a candidate under the fixed
rules. Campaign records distinguish these decisions from initial
configuration, file transfers, restarts, code edits, and automated retries.
Each has its own count and timestamp; absence of a record remains unknown.
The intended steady-state operation has zero routine human scientific
decisions after initialization.

Valid attempts, no-promotion patience, abstention, nonproductive portfolios,
integrity checks, and total resource limits determine stopping.
Infrastructure-invalid attempts follow bounded recovery rules and retain
their cost even when excluded from the scientific-attempt counter.
The observed campaign stops after five valid rounds and three successive
rounds without promotion. The invalid attempt remains separately counted;
later reused-evidence validations do not increment the formal round count.
Local Analyzer/Planner distillation is a separate research-role experiment,
with shadow qualification preceding any takeover. The main method can use
fixed strong research models throughout.
